
\section{Is There a Difference?}

Editorial article in "Folkebladet," August 4, 1897.

There is no little lack of clarity among many as to whether there is any difference between the Free Church and the old "church bodies." We will try to shed a little light on this matter.

By "the Free Church" is meant the free cooperation of free congregations. By "church bodies" is meant a centralized common government which uses congregations and pastors in its service for the advancement of its own aims. There proves to be no little difference between these two things, when it comes down to it.

It is the same difference as between many small landowners and farmers, each of whom builds his own home and cultivates his own piece of land, on the one side, and, on the other, one great and mighty estate-owner with many tenant farmers under him.

Suppose one were to try these two methods, many small farms and one great estate, side by side under conditions as nearly equal as possible: the same kind of soil, the same kind of climate, the same kind of products and crops; where would the land then be best used, and where would the people’s well-being and prosperity thrive most?

Take one "large farm" of 10,000 acres and set beside it one hundred "small farms" of 100 acres each; where would there be the happier population, the better farming, and the greater prosperity? Where would there be the most real yield? Or would there be no difference?

It is true that in the statistics it looks rather impressive when the large farmer comes with his heavy figures, his thousands of bushels of every kind, while the small farmers have only some hundreds each. But when the reckoning is carried a little further, it comes to matter, after all, whether the figures have counted rightly, and whether glad hearts and happy homes are also measured by the bushel.

Experience shows that there is no better way to build up a country than that each man own his home and his land and use and cultivate it in such a way as is serviceable for him. Experience also shows that the people thrive best, and peace dwells most securely in the community, where one man stands about as high as another, and all are equally free to do with their own as they will.

They cooperate none the less. With each man’s prosperity the community’s grows; with each plot that is made more beautiful, the whole community becomes fairer. Can this be avoided, or is it not self-evident?

Think carefully upon this picture, and see whether it does not fit the Free Church well enough. All that is meant by the Free Church is the independence and freedom of the congregations, so that each one for itself may become a true spiritual congregation, which seeks to do its work as well as possible, because in that are God’s good pleasure and its own true good.

That does not hinder them from cooperating. It does not hinder them from learning from one another, helping one another, supporting one another, encouraging one another, living together and working together in love.

On the contrary, it is inevitable. If there is spiritual life in one congregation, this will of itself have its effect upon many other congregations.

In English we have two words which roughly express the difference between the Free Church and church bodies. These are "Co-operation" and "Concentration." By "Co-operation" is meant that many forces work together because they have common interests which they must look after; by "Concentration" is meant that there is a central power which directs the labor of all, because only thus does one attain the greatest possible unity in the work.

These two methods are very different from one another; and in the ecclesiastical sphere they lead to very different results. That the Free Church is both the biblical and the Lutheran way, and therefore also the one which brings the best results, upon that point the learned ought not to be in disagreement.

Georg Sverdrup.
